\documentclass[11pt]{article}

\usepackage[margin=1in]{geometry}
\usepackage[T1]{fontenc}
\usepackage{array}
\usepackage{longtable}
\usepackage{hyperref}

\title{Latency and Reliability Validation of Three-Channel Audio--Tactile Stimulus Administration}
\author{Peripersonal Space Toolkit Internal Validation}
\date{Living report, updated 2026-06-13}

\newcommand{\implemented}{\textsc{implemented}}
\newcommand{\reviewrequired}{\textsc{review required}}
\newcommand{\proven}{\textsc{proven}}
\newcommand{\deferred}{\textsc{deferred}}

\begin{document}
\maketitle

\begin{abstract}
This report records the validation strategy for the Peripersonal Space Toolkit
timing and data-recording stack. The immediate technical question is whether a
single three-channel WAV file can be administered through the Komplete Audio 6
so that channel 1 drives the left headphone path, channel 2 drives the right
headphone path, and channel 3 drives the Woojer/tactile path, with all three
channels remaining synchronized. The second question is whether the software
measurement chain, including local event logs, Lab Streaming Layer (LSL)
markers, physical loopback recordings, LabRecorder/XDF records, and emulated
mouse-click responses, measures those events with enough precision and
reliability to support experiment interpretation. This is an internal
validation report. It does not replace participant-facing documentation; the
Woojer device is not physically in the current loop, so the report claims
electrical channel-3 tactile-drive timing rather than mechanical vibration
onset.
\end{abstract}

\section{Report Scope}

This document records the internal validation evidence for the current
three-channel Komplete Audio 6 setup. The main scientific claim is scoped to
electrical stimulus administration and data recording: one dummy three-channel
WAV is split to three physical outputs, the output channels remain synchronized
within the measured tolerances, local logs and LSL/XDF preserve the event
records, and direct electrical loopback provides the physical timing reference.

The Woojer transducer is not physically included in this validation loop.
Accordingly, the report treats channel 3 as the electrical tactile-drive path.
Mechanical Woojer vibration onset is a separate later measurement that requires
the Woojer plus an external vibration/contact sensor.

\section{Quantitative Reporting Standard}

The scientific end report must report timing values as distributions, not only
as pass/fail statements. For each accepted physical timing dataset, the report
should include the number of detected pulses, mean latency, standard deviation
(SD), median, p95, minimum, and maximum for each channel. Inter-channel
synchronization should be reported from paired pulse events as mean absolute
left/right skew $\pm$ SD and mean absolute tactile-vs-audio skew $\pm$ SD,
with the paired-pulse count stated explicitly. These are the primary numerical
answers to whether one three-channel WAV remains synchronized across headphone
left, headphone right, and tactile-drive output.

The final report must also compare the reliability of the active recording
layers separately: local logs, LSL/XDF, and physical electrical loopback. For
each layer, report marker or pulse detection rate, missing/duplicate events,
timing error mean $\pm$ SD where a reference exists, and the layer's
interpretation boundary. LSL is the rich software reconstruction stream, local
logs are the primary software record, and direct electrical loopback is the
ground-truth physical timing layer for output synchronization. WASAPI loopback
was evaluated and excluded from the core ASIO validation strategy because ASIO
multichannel playback can bypass the Windows endpoint that WASAPI records.

\section{Plain-Language Summary}

The purpose of this validation is simple: before trusting real experiment
stimuli, we need to know whether the computer and audio interface can send one
three-lane audio file to the right places at the right time. In this setup,
lane 1 should go to the left headphone, lane 2 should go to the right
headphone, and lane 3 should go to the tactile device path that can feed the
Woojer. We also need to know whether the recording system can later tell us
what really happened.

To test this, we made an intentionally simple three-channel WAV file. Instead
of using realistic looming sounds, it contains five easy-to-recognize pulse
events. Each channel has a slightly different pulse shape. That means the
analysis can tell whether a channel arrived on the correct input, whether a
channel was routed incorrectly, whether one channel disappeared, or whether the
signal was mixed together.

The accepted finding so far is that the core administration path is feasible:
the Komplete Audio 6 ASIO driver can open and play a stable three-channel
stream, and the validation tooling can test whether a single multichannel WAV
is preserved as independent left, right, and tactile-drive outputs.

For publication-style reporting, this document separates accepted evidence from
setup and acquisition checks. Numerical inter-channel latency and skew should
only be reported from a final matched-gain direct-loopback run in the functional
route state, where all three channels pass the same detection, clipping, and
signal-quality criteria. Earlier setup checks are retained locally for
engineering audit, but they are not used as publication-facing estimates.

The LSL data-recording side worked well after installing the optional LSL
dependency. A dummy marker stream delivered all 5 expected pulse markers to an
external probe. The rich experiment-style LSL stream,
\texttt{PPSMarkersV2}, delivered all markers in the tested bursts, including
event type, event code, trial identity, sample index, and timestamp-quality
labels. The numeric trigger-code stream also worked when the test allowed a
short warmup and sent markers at a realistic pace. A first extremely short
dual-stream burst was missed by the external probes because the sender opened,
emitted, and closed too quickly. The practical lesson is that LSL outlets
should be created before the first important marker, kept alive for the run,
and not used as one-shot throwaway streams.

An independent LabRecorder/XDF stress test also passed. LabRecorderCLI captured
16/16 rich \texttt{PPSMarkersV2} samples and 16/16 numeric
\texttt{PPSTriggerCodes} samples, with zero missing event IDs, zero extra event
IDs, zero duplicate event IDs, zero field mismatches, and matching numeric code
counts. The XDF timestamps differed from the local marker mirror by about
-0.005 ms on this run. This validates the external software recording path; it
does not by itself prove physical signal arrival at the audio outputs.

The actual-condition recording-layer alignment run now joins the software and
physical evidence in one test. One prepared 88 s Segment 5/6 block was played
through the Komplete ASIO route while the runner wrote its local event/LSL
marker mirror, wrote an optional local digital audio evidence WAV from the
mixed output callback, and captured a validation-only physical electrical
loopback. The comparison passed. Physical-minus-digital latency was
33.462 $\pm$ 0.013 ms, the left/right skew was -0.023 ms, the
tactile-vs-audio skew was 0.011 ms, the digital audio evidence WAV had zero
dropped buffers and no clipping, and the internal marker mirror contained all
146 runtime event IDs with no missing or duplicate markers. Callback-derived
LSL timestamps were evaluated against the audio-sample timeline after anchoring
events to \texttt{audio\_sample\_zero}; p95 timestamp error was below
0.001 microseconds in the local reconstruction.

WASAPI loopback did not record usable data during ASIO multichannel dummy
playback in the tested run. This is expected for this setup: the Native
Instruments ASIO stream can bypass the Windows WASAPI endpoint. For the
Komplete ASIO three-channel route, WASAPI is therefore removed from the core
validation strategy. The direct electrical loopback capture is the physical
record of what left the interface.

Finally, the Woojer itself is not physically in the current loop. These tests
measure the electrical signal sent toward the tactile path. When the strap is
added later, mechanical vibration onset should be measured with an external
vibration sensor, accelerometer, or contact microphone.

\section{Current Findings Snapshot}

The main validation goal is still the same: determine whether one three-channel
stimulus file can be routed through the Komplete Audio 6 with synchronized
left-audio, right-audio, and tactile-drive channels, and whether the experiment
can record what happened with enough precision to trust later analysis.

The strongest positive finding is on the software and recording side. The
Komplete ASIO endpoint can sustain the required multichannel stream, the local
event log and rich/numeric LSL marker streams reconstruct the same actual event
set in the tested stress runs, and LabRecorder/XDF preserves those LSL streams
without missing or duplicating markers. The current evidence audit lists these
software/data-recording requirements as proven.

The response-click software path is also now supported at three levels. The
synthetic timing-event stress linked 50/50 mouse clicks to 50/50 response-marker
events. The session-controller stress linked 25/25 clicks to 25/25
callback-derived markers, with marker-minus-mouse delay mean 8.150 ms, SD 0.081
ms, and median 8.105 ms for an intended 8 ms marker. The visible legacy Tk
runner OS-click stress sent 10 armed OS clicks into the click target during
fake-audio playback and recovered 10/10 in-target mouse clicks plus 10/10
linked response markers, with marker-minus-mouse delay mean 8.553 ms, SD 0.125
ms, and median 8.541 ms. This validates OS-click delivery into the current
visible runner path without touching the audio hardware.

A realtime one-block runner stress now verifies the immediate data-output
contract. The harness prepared one Segment 5/6-style participant block, ran it
on the runner timeline with a deterministic fake audio engine, and injected
jittered clicks after tactile onsets. The current run recovered 5/5 tactile
onsets, 5/5 mouse clicks, 5/5 response markers, and 5/5 trial-end markers. It
wrote \texttt{events.csv}, a pyxdf-loadable \texttt{events.xdf},
\texttt{lsl\_markers.csv}, \texttt{trigger\_dictionary.json}, timing QC, and a
clean \texttt{analysis\_ready\_trials.csv}. The five simulated response rows
had mean RT 223.215 ms, SD 30.936 ms, and 5/5 hits. Mouse-to-response-marker
delay was 8.081 ms mean, SD 0.014 ms.

The next-phase actual-condition one-block recording-layer run also passed. A
real prepared Pfeiffer-style Segment 5/6 block was administered through the
Komplete ASIO route as one continuous 3-channel block WAV with conservative
validation gains (\texttt{audio\_gain = 0.0005},
\texttt{tactile\_gain = 0.02}). The session used
\texttt{participant\_block\_wavs}, contained exactly one block, emitted 22/22
trial starts, 12/12 looming onsets, 20/20 tactile onsets, 22/22 response-window
onsets, and 22/22 trial ends, and used no timing-anchor fallback. The internal
XDF loaded with 146 samples, the LSL marker mirror contained 146 rows, the
analysis-ready trial table contained 20 response rows, and the optional digital
audio evidence WAV was written with zero dropped buffers and no clipping. The
matching direct electrical capture was unclipped. Recording-layer comparison
estimated physical-minus-digital latency at 33.462 $\pm$ 0.013 ms
(median 33.469 ms, p95 33.469 ms, range 33.447--33.469 ms), left/right skew at
-0.023 ms, and tactile-minus-audio-mean skew at 0.011 ms. The internal
callback-derived LSL/event timestamp error relative to the sample timeline was
0.00000019 $\pm$ 0.00000035 ms, with p95 below 0.000001 ms. The physical
tactile-channel response-marker trace recovered 20/20 markers; after fitting
the common recording offset, residual jitter was 0.000 ms.

The physical channel-latency question is now supported by a clean
channel-scaled dummy validation run. A safe route sweep confirmed the intended
one-to-one physical identities: output 1 to input 1, output 2 to input 2, and
output 3 to input 3. The accepted simultaneous three-channel dummy WAV direct
loopback then recovered all five planned pulses on all three channels with no
low-signal or clipping flags. Mean electrical latency was 33.605 ms on channel
1, 33.583 ms on channel 2, and 33.605 ms on channel 3; all three channel SDs
were 0.000 ms in the accepted five-pulse run. Paired left/right electrical skew
was 0.023 ms mean, SD 0.000 ms. Paired tactile-vs-audio electrical skew was
0.011 ms mean, SD 0.000 ms. This validates the electrical synchronization of
the one-file three-channel route under the tested functional setup.

The physical boundary is therefore explicit: the accepted result is electrical
channel-3 drive timing at the Komplete route. Mechanical Woojer vibration
timing is a later hardware phase because the Woojer is not physically included
in the current validation loop.

\section{Validation Questions}

The validation program is designed around four questions.

\begin{enumerate}
  \item Can the current software and Komplete Audio 6 setup play a single
  three-channel WAV as three independent physical outputs?
  \item Are channels 1 and 2, used for the Sennheiser headphone paths, and
  channel 3, used for the Woojer/tactile path, electrically synchronized within
  the required tolerance?
  \item Can the software recording chain measure expected stimulus timing,
  trigger timing, LSL markers, direct loopback traces, and mouse-click response
  markers without missing, duplicating, or systematically misplacing events?
  \item Can the full experiment be reconstructed from callback-derived local
  logs and LSL marker mirrors, while retaining an optional local audio evidence
  WAV as a data-heavy safety copy? Physical loopback is the validation reference
  for this question, not a required participant-run layer.
\end{enumerate}

\section{System Under Test}

The system under test is the installed PPS Toolkit repository, the Windows
audio stack, and the lab hardware route. Validation scripts are intentionally
kept outside the package in \texttt{validation\_protocols/}. They exercise the
existing runtime from the outside rather than adding a calibration interface or
participant-facing feature.

The current operational runner surface is native Focus Mode /
\texttt{PPSExperimentRunner.exe}. The older visible Tk runner was validated as
historical regression evidence for the shared timing core, LSL streams,
callback timestamps, response-marker ordering, and loopback requirements, but
it is no longer an operator launch path.

For next-phase validation, the accepted unit of runner evidence is exactly one
actual Segment 5/6 prepared experimental block. Dummy pulse acquisitions and
fake-audio runner stresses remain useful preflights, but they are not counted
as formal actual-condition runner evidence. The offline validator
\texttt{validate\_one\_block\_actual\_condition\_run.py} audits a completed
session folder and fails runs that are not one-block
\texttt{participant\_block\_wavs} sessions, are missing analysis/XDF/LSL
outputs, use timing fallbacks, or are sourced from dummy/fake validation
fixtures. The validation-only harness
\texttt{run\_one\_block\_actual\_condition\_validation.py} can repeat this
test with full-duplex ASIO direct capture and deterministic simulated clicks;
it is not a public participant workflow.

\begin{longtable}{|p{0.26\textwidth}|p{0.66\textwidth}|}
\hline
\textbf{Component} & \textbf{Role in validation} \\
\hline
Komplete Audio 6 MK2 & Multichannel ASIO output device. The intended route is
one synchronized endpoint rather than separate stereo Windows endpoints. \\
\hline
Three-channel WAV & The dummy stimulus source. Channel 1 is left headphone,
channel 2 is right headphone, and channel 3 is tactile/Woojer drive. \\
\hline
Direct electrical loopback & Primary ground truth for interface-level channel
latency, skew, crosstalk, downmixing, clipping, and channel loss. \\
\hline
Local audio evidence WAV & Optional experiment-time software safety copy. It is
written from the runner's mixed output callback after routing, gain, clipping
limits, and tactile-channel response-marker injection. It proves the buffers the
runner attempted to send, while physical loopback proves electrical arrival. \\
\hline
WASAPI loopback & Evaluated and excluded from the core acceptance strategy for
this ASIO multichannel route. It records Windows endpoints, while the Native
Instruments ASIO stream can bypass that path. \\
\hline
LSL rich stream & The runner's reconstructable marker stream
\texttt{PPSMarkersV2}. This stream should contain enough event identity and
timing metadata to decode the experiment from LSL alone. \\
\hline
LSL numeric stream & The runner's \texttt{PPSTriggerCodes} stream for
trigger-code pipelines. Reconstruction depends on the matching
\texttt{trigger\_dictionary.json}. \\
\hline
Local event logs & \texttt{events.csv}, \texttt{events.xdf},
\texttt{lsl\_markers.csv}, timing QC files, and session analysis outputs.
These are the non-LSL record. \\
\hline
Mouse-click emulation & Internal stress input for response timing. The goal is
to compare immediate software click logging with tactile-channel physical
response markers. \\
\hline
Woojer mechanical validation & Deferred later stage. The Woojer is not
physically in the current loop, so the present claim is electrical channel-3
tactile-drive timing. Mechanical onset requires a contact microphone,
accelerometer, or other external vibration sensor. \\
\hline
\end{longtable}

\section{Software And Driver Requirements}

The validation depends on a specific Windows software stack as well as the
physical audio route. The tracked install checklist is
\texttt{docs/WINDOWS\_PC\_SOFTWARE\_REQUIREMENTS.md}. The current lab install
uses 64-bit Python 3.12, the toolkit installed with
\texttt{.[gui,web,lsl,validation]}, the official Native Instruments
\texttt{Komplete Audio ASIO Driver}, MiKTeX/\texttt{pdflatex} for rebuilding
this report, and LabRecorder for independent LSL/XDF recording.

The dependency audit script
\texttt{validation\_protocols/scripts/audit\_pc\_software\_requirements.py}
checks Python packages, external tools, the ASIO registry, and whether
\texttt{sounddevice} can see the Komplete ASIO endpoint as a full-duplex
multichannel device. On 2026-06-12, the audit found the runtime stack, GUI/web
extras, LSL package, \texttt{pyxdf}, MiKTeX, Git, LabRecorder, and Komplete
ASIO endpoint present on the validation PC. The only downloaded external
recording tool was the official LabRecorder Windows release zip, stored under
ignored \texttt{local\_data/software\_installers/} and extracted under ignored
\texttt{local\_data/software\_For-AI/engineering/tooling/}; third-party binaries are not committed to
the repository.

The Native Instruments driver was not replaced or reinstalled by the validation
scripts. Driver installation remains a manual vendor-controlled step via Native
Instruments or Native Access. This is deliberate: the publication timing route
requires the native NI ASIO endpoint, while generic wrappers or split Windows
stereo endpoints remain diagnostic only.

\section{Measurement Definitions}

The validation separates latency sources that are often conflated.

\begin{longtable}{|p{0.22\textwidth}|p{0.44\textwidth}|p{0.22\textwidth}|}
\hline
\textbf{Quantity} & \textbf{Definition} & \textbf{Primary artifact} \\
\hline
Planned sample timing & Expected pulse or trial-event sample index in the WAV
or block CSV. & \texttt{planned\_pulses.csv},
\texttt{block\_XX.csv} \\
\hline
Callback timestamp & Timestamp assigned when a scheduled event crosses the
audio output callback buffer. Prefer DAC-time-derived timestamps when exposed
by PortAudio/sounddevice. & \texttt{events.csv},
\texttt{lsl\_markers.csv} \\
\hline
Electrical output latency & Delay between planned output sample and looped-back
input detection. This includes interface output/input buffering and physical
analog route. & direct loopback WAV and comparison report \\
\hline
Digital-vs-physical offset & Difference between the local digital audio evidence
WAV and the direct physical loopback WAV for the same output buffer timeline. &
\texttt{recording\_layer\_alignment\_report.json} \\
\hline
Channel skew & Difference between median detected latencies across channels.
This is the key measure for left/right and audio/tactile synchrony. & dummy
pulse comparison report \\
\hline
LSL arrival behavior & Difference between marker timestamp and probe arrival
clock. This measures data-collection behavior, not physical signal arrival. &
\texttt{lsl\_marker\_probe.csv} \\
\hline
Mouse-to-marker delay & Delay from immediate software mouse-click event to the
response marker entering the tactile output path. & \texttt{timing\_qc.csv},
loopback recording \\
\hline
Woojer mechanical onset & Physical vibration onset after the Woojer receives
channel 3. & \deferred{} until the Woojer is physically added with an external
sensor \\
\hline
\end{longtable}

\section{Dummy Three-Channel Validation}

The first-run validation uses a deliberately simple dummy stimulus rather than
a real experiment stimulus. The script
\texttt{validation\_protocols/scripts/make\_dummy\_pulse\_stimulus.py}
generates a three-channel WAV with five rectangular coded pulse events. The
inter-pulse intervals vary so the trace is easy to recognize during analysis.
The channels share nominal pulse onsets, but each channel has a different pulse
shape. This makes misrouting, downmixing, crosstalk, or channel-3 loss
detectable.

\subsection{Generated Artifacts}

Each dummy run produces:

\begin{itemize}
  \item \texttt{dummy\_3ch\_pulses.wav}, the source three-channel stimulus.
  \item \texttt{planned\_pulses.csv} and \texttt{planned\_pulses.json}, the
  expected event grid by channel and pulse index.
  \item \texttt{dummy\_pulse\_manifest.json}, the stimulus and route metadata.
  \item \texttt{direct\_loopback\_capture.wav}, when direct loopback is armed.
  \item \texttt{lsl\_emitted\_markers.csv}, when dummy LSL emission is enabled.
  \item \texttt{dummy\_pulse\_comparison\_report.json} and
  \texttt{dummy\_pulse\_comparison\_report.md}, after comparison.
\end{itemize}

\subsection{Primary Commands}

\begin{verbatim}
python .\For-AI/engineering/validation\scripts\make_dummy_pulse_stimulus.py

python .\For-AI/engineering/validation\scripts\run_dummy_pulse_latency.py ^
  --device-query Komplete --record-asio-loopback

python .\For-AI/engineering/validation\scripts\run_dummy_pulse_latency.py ^
  --device-query Komplete --record-asio-loopback --emit-lsl

python .\For-AI/engineering/validation\scripts\compare_dummy_pulse_recordings.py ^
  --run-dir artifacts\validation_runs\dummy_pulse_YYYYMMDD_HHMMSS
\end{verbatim}

\subsection{Acceptance Criteria}

\begin{longtable}{|p{0.42\textwidth}|p{0.18\textwidth}|p{0.28\textwidth}|}
\hline
\textbf{Criterion} & \textbf{Threshold} & \textbf{Interpretation} \\
\hline
Three-channel capture is present & at least 3 channels & Confirms the capture
contains all routed outputs. \\
\hline
Channel identity map & 1 to 1, 2 to 2, 3 to 3 & Detects unintended channel
assignment, downmixing, or channel-3 loss. \\
\hline
Pulse detection rate & at least 95\% per channel & Confirms that all expected
pulses are recoverable at the accepted acquisition level. \\
\hline
Left/right electrical skew & at most 1 ms & Required for synchronized
headphone delivery. \\
\hline
Tactile/audio electrical skew & at most 2 ms & Required for channel-3 timing
to be interpreted relative to channels 1/2. \\
\hline
Residual jitter & p95 at most 2 ms, max at most 5 ms & Confirms stable capture
latency across repeated pulses. \\
\hline
Signal quality & no clipping, no low-signal failure & Confirms usable
recording amplitude. \\
\hline
LSL marker count, when enabled & all expected pulse IDs, no duplicates &
Confirms marker reliability for the dummy timing grid. \\
\hline
\end{longtable}

\section{Experiment-Runner LSL Validation}

The experiment runner should publish a rich LSL marker stream named
\texttt{PPSMarkersV2}. Its current channel order is:

\begin{verbatim}
marker_version, event_id, event_type, event_code, trigger_key,
session_id, participant_id, block_index, trial_uid, sample_index,
timestamp_quality, payload_json
\end{verbatim}

The runner may also publish \texttt{PPSTriggerCodes}, a numeric
\texttt{int32} trigger stream. For EEG-style systems that record only numeric
triggers, event reconstruction depends on the session's deterministic
\texttt{trigger\_dictionary.json}. For full LSL reconstruction, the rich stream
must contain the event identity, trial identity, code, trigger key, sample
index, timestamp-quality label, and payload JSON.

The external probe
\texttt{validation\_protocols/scripts/lsl\_marker\_probe.py} now records
\texttt{PPSMarkersV2} by default and normalizes both rich v2 samples and older
compact three-field marker samples into a comparable CSV. The acceptance check
uses \texttt{reconcile\_lsl\_with\_local\_events.py} to compare the external
probe CSV against \texttt{events.csv}, \texttt{lsl\_markers.csv}, and, for the
numeric stream, the expected trigger-code counts. Missing event IDs, duplicate
event IDs, trigger-key mismatches, unexpected timestamp-quality changes, and
numeric code-count mismatches require review.

For an external-recorder validation, the script
\texttt{validation\_protocols/scripts/run\_labrecorder\_lsl\_xdf\_stress.py}
creates the same rich and numeric PPS LSL streams, records them with
LabRecorderCLI, loads the resulting XDF with \texttt{pyxdf}, and compares the
recorded samples against the local \texttt{lsl\_markers.csv} mirror and
\texttt{trigger\_dictionary.json}. This tests whether a standard external LSL
recorder preserves the data needed for experiment reconstruction.

\section{Mouse-Click and Response-Marker Validation}

Response validation has two complementary paths. First, a mouse-click event is
logged immediately by the software. Second, an associated low-gain
response-marker pulse is sent to the tactile channel so a physical loopback
recording can recover the response event if local logs or LSL are incomplete.
The emulated-click protocol does not estimate human reaction time. It tests
whether software click capture, LSL click markers, local event logs, response
marker scheduling, and physical tactile-channel marker output remain linked
and stable under deterministic input.

The response strategy comparison separates three clocks that are easy to
confuse: the local mouse-click log, the rich LSL sample timestamp, and the time
an external LSL probe receives the sample. In the current 50-click stress run,
local click-to-response-marker delay was 8.000 ms mean, SD effectively 0.000
ms, and median 8.000 ms. Rich LSL mouse-click sample timestamps followed local
mouse times by 0.081 ms mean, SD 0.056 ms, and median 0.067 ms, while external
probe arrival followed by 0.421 ms mean, SD 0.223 ms, and median 0.339 ms. The
callback-derived response-marker sample timestamp occurred about 8.005 ms after
the local mouse event, matching the planned 8 ms software delay.
Response-marker LSL arrival happened earlier than the marker timestamp because
the marker is pushed with an explicit future DAC/sample time; arrival time
should therefore not be interpreted as physical marker onset.

A second runner-layer stress test now exercises
\texttt{SessionRunnerController.log\_click()} during active playback with a
deterministic fake audio engine. This validates the actual session-controller
response path without sending OS mouse clicks or touching hardware. The current
runner order is deliberately conservative: log the click locally, trigger the
tactile response marker, and only then push the already logged mouse marker to
LSL. In the deterministic runner-layer stress, this logged 25/25 mouse clicks
and 25/25 linked response markers, all with
\texttt{dac\_time\_sample\_exact} marker timestamps; marker-minus-mouse delay
was 8.150 ms mean, SD 0.081 ms, median 8.105 ms, and maximum 8.317 ms for an
intended 8 ms marker.

A third stress test opens the retired visible Tk runner window and sends armed
OS mouse clicks to the click target while a deterministic fake audio engine
simulates active playback. This historical regression test originally exposed a
Windows DPI-scaling coordinate mismatch in the click-target hit test; the runner
now compares pynput mouse coordinates against DPI-adjusted Tk widget bounds.
After that fix, the visible-runner stress recorded 10/10 in-target OS clicks
during playback and 10/10 linked response markers. Marker-minus-mouse delay was
8.553 ms mean, SD 0.125 ms, median 8.541 ms, and p95 8.762 ms. This validates
GUI-level click
delivery in the current runner without using the Komplete Audio 6 outputs.
The refreshed screenshot-enabled pass recovered 3/3 in-target OS clicks and
3/3 linked response markers, with marker-minus-mouse delay 8.335 ms mean, SD
0.086 ms, and a visible runner screenshot showing the Recording row for LSL,
XDF, analysis CSV, and audio-evidence choices.

A fourth runner stress validates the complete one-block software output path.
\texttt{run\_one\_block\_trial\_runner\_realtime\_stress.py} creates a
Segment 5/6-style fixture, prepares one participant block, runs
\texttt{SessionRunnerController} in realtime with a fake audio engine, injects
jittered clicks after tactile onsets, and verifies that the resulting files are
immediately analyzable. The current run recovered 5/5 tactile onsets, 5/5 mouse
clicks, 5/5 response markers, and 5/5 trial-end markers, with no duplicate
event IDs. The internal \texttt{events.xdf} loaded with \texttt{pyxdf} and
contained 43 samples. The copied \texttt{analysis\_ready\_trials.csv} contained
5/5 hit rows with RT mean 223.215 ms and SD 30.936 ms.

\section{Current Evidence}

This section records validation evidence collected on 2026-06-12. The hardware
tests were run on the local Windows machine with the Komplete Audio ASIO Driver
visible as device 31, exposing 6 inputs and 6 outputs at 44100 Hz. The accepted
hardware finding is that the Komplete ASIO endpoint can sustain synchronized
three-channel output at the intended sample rate, block size, and latency
settings. The accepted dummy calibration provides the clean electrical latency
baseline, and the accepted actual-condition one-block run verifies the same
runner output contract on a real prepared experimental block.

\begin{longtable}{|p{0.36\textwidth}|p{0.18\textwidth}|p{0.34\textwidth}|}
\hline
\textbf{Item} & \textbf{Status} & \textbf{Evidence} \\
\hline
Internal validation folder separated from package runtime & \implemented{} &
\texttt{validation\_protocols/} contains protocols, scripts, templates, and
this report outside \texttt{src/}. \\
\hline
Dummy coded three-channel WAV generation & \implemented{} &
\texttt{make\_dummy\_pulse\_stimulus.py}; covered by
\texttt{tests/test\_validation\_protocols.py}. \\
\hline
Direct-loopback dummy runner & \implemented{} &
\texttt{run\_dummy\_pulse\_latency.py} can generate, play, capture, and record
run metadata when hardware is connected. \\
\hline
Dummy capture comparison & \implemented{} &
\texttt{compare\_dummy\_pulse\_recordings.py} detects channel identity,
latency, skew, jitter, clipping, low signal, and pass/fail criteria. \\
\hline
Synthetic comparison tests & \implemented{} &
Unit tests recover known latency and detect misrouted channels. \\
\hline
Komplete direct loopback hardware result & \proven{} & The channel-scaled
functional route sweep and simultaneous dummy direct-loopback run accepted all
three channels under the same detection, signal-quality, and no-clipping
criteria. Electrical latency and skew values are now reportable for the tested
functional route. \\
\hline
WASAPI boundary result & \proven{} & WASAPI loopback was requested during ASIO
dummy playback and initialized, but saved no data. This establishes that WASAPI
is not the acceptance path for this Komplete ASIO three-channel route; the ASIO
stream can bypass the Windows endpoint that WASAPI records. \\
\hline
LSL probe of \texttt{PPSMarkersV2} & \implemented{} & A deterministic
\texttt{TimingEventHub} burst emitted 12 rich markers; the external probe
received all 12 with no duplicate event IDs and preserved event type, event
code, trigger key, trial UID, sample index, and timestamp-quality fields. \\
\hline
LSL/local reconciliation & \implemented{} & The mouse/response timing stress
run reconciled 104/104 local actual events against 104/104 rich
\texttt{PPSMarkersV2} samples with zero missing, extra, duplicate, or
field-mismatched event IDs; the numeric \texttt{PPSTriggerCodes} stream also
matched all expected trigger-code counts. \\
\hline
LabRecorder/XDF preservation & \implemented{} & LabRecorderCLI captured 16/16
rich \texttt{PPSMarkersV2} samples and 16/16 numeric
\texttt{PPSTriggerCodes} samples. Comparison against local marker mirrors found
zero missing, extra, or duplicate event IDs, zero metadata field mismatches,
and matching numeric trigger-code counts. \\
\hline
Mouse/response timing stress result & \implemented{} & Internal software
harness logged 50/50 mouse clicks and 50/50 linked response markers, with
8.000 ms monotonic marker-minus-mouse timing and 104/104 rich plus 104/104
numeric LSL samples captured by external probes. \\
\hline
Response timing strategy comparison & \implemented{} & The existing
mouse/response stress artifacts were compared across local logs and rich LSL
probe timing. Local click-to-marker delay was 8.000 ms mean with negligible SD;
rich LSL mouse-click sample timestamps trailed local clicks by 0.081 ms mean,
SD 0.056 ms, while external LSL arrival trailed by 0.421 ms mean, SD 0.223 ms.
\\
\hline
Session runner click-path stress & \implemented{} & Deterministic
\texttt{SessionRunnerController.log\_click()} stress with a fake audio engine
logged 25/25 clicks and 25/25 linked response markers during active playback.
After deferring LSL publishing until after marker trigger, marker-minus-mouse
delay was 8.150 ms mean, SD 0.081 ms, median 8.105 ms, and maximum 8.317 ms for
an intended 8 ms marker. \\
\hline
Visible legacy-runner OS-click stress & \implemented{} &
\texttt{run\_visible\_runner\_os\_click\_stress.py} opened the current Tk
runner with a fake audio engine and sent armed OS clicks to the target. The
passing run recovered 10/10 in-target clicks and 10/10 linked response markers;
marker-minus-mouse delay was 8.553 ms mean, SD 0.125 ms, median 8.541 ms. \\
\hline
One-block realtime runner output stress & \implemented{} &
\texttt{run\_one\_block\_trial\_runner\_realtime\_stress.py} prepared one
Segment 5/6-style block, ran it on the runner timeline, injected jittered
post-tactile clicks, and verified immediate local outputs. The current run
recovered 5/5 tactile onsets, mouse clicks, response markers, and trial-end
markers; \texttt{events.xdf} loaded with \texttt{pyxdf}; and
\texttt{analysis\_ready\_trials.csv} contained 5/5 hit rows. \\
\hline
Actual-condition one-block runner output & \proven{} &
\texttt{run\_one\_block\_actual\_condition\_validation.py} administered one
real prepared Pfeiffer-style Segment 5/6 block through Komplete ASIO with
full-duplex direct capture. The accepted run passed the one-block validator:
22/22 trial starts, 12/12 looming onsets, 20/20 tactile onsets, 22/22 response
windows, 22/22 trial ends, 146 loadable XDF samples, 146 LSL mirror rows, and
20 analysis-ready response rows. \\
\hline
Actual-condition direct loopback comparison & \proven{} &
\texttt{compare\_actual\_block\_loopback.py} compared the actual block WAV
against the matching 3-channel direct capture. The capture was unclipped.
Summary-eligible correlated segments yielded left/right skew -0.023 $\pm$
0.000 ms and tactile-minus-audio-mean skew -0.079 $\pm$ 0.000 ms. Absolute
capture alignment includes capture lead-in and is not interpreted as pure
hardware latency. \\
\hline
Response-marker loopback comparison pipeline & \implemented{} &
\texttt{compare\_response\_marker\_loopback.py} reads logged
\texttt{response\_marker\_start} sample indices and tactile-channel recordings,
fits the recording/hardware offset, and reports detection rate plus residual
jitter. Synthetic regression coverage verifies exact recovery when marker
pulses are present in a known loopback trace. \\
\hline
Woojer mechanical vibration onset & \deferred{} & The Woojer is not physically
in the current validation loop. The report therefore claims electrical
channel-3 tactile-drive timing only. \\
\hline
\end{longtable}

\subsection{Artifact Index}

\begin{longtable}{|p{0.30\textwidth}|p{0.58\textwidth}|}
\hline
\textbf{Artifact} & \textbf{Validation use} \\
\hline
\texttt{artifacts/validation\_runs/dummy\_output\_route\_sweep\_*}
& Internal direct-loopback route identity and level-check artifacts. These
setup-level records are used to verify acquisition readiness and refine the
accepted baseline protocol; they are not treated here as publication-grade
latency/skew baselines. \\
\hline
\texttt{artifacts/validation\_runs/dummy\_pulse\_lsl\_20260612\_131637} and
\texttt{artifacts/validation\_runs/lsl\_probe\_dummy\_pulse\_20260612\_131637}
& Dummy pulse run with LSL emission plus external probe; 5/5 dummy markers
received. \\
\hline
\texttt{artifacts/validation\_runs/lsl\_v2\_burst\_20260612\_131948} and
\texttt{artifacts/validation\_runs/lsl\_probe\_v2\_burst\_20260612\_131948}
& Rich \texttt{PPSMarkersV2} deterministic burst and external probe; 12/12
markers received. \\
\hline
\texttt{artifacts/validation\_runs/lsl\_dual\_stream\_burst\_20260612\_132503}
& First simultaneous rich/numeric LSL probe attempt; both streams resolved but
0 samples were captured because the synthetic outlet emitted and closed too
quickly. This is retained as evidence that one-shot LSL bursts are not a robust
strategy. \\
\hline
\texttt{artifacts/validation\_runs/lsl\_dual\_stream\_burst\_paced\_20260612\_132607}
with paired rich/numeric probe folders & Paced simultaneous LSL test after
adding outlet warmup and marker spacing. Both \texttt{PPSMarkersV2} and
\texttt{PPSTriggerCodes} were received 8/8. \\
\hline
\texttt{artifacts/validation\_runs/labrecorder\_lsl\_xdf\_current}
& External LabRecorderCLI/XDF stress run. LabRecorder captured 16/16 rich
\texttt{PPSMarkersV2} samples and 16/16 numeric \texttt{PPSTriggerCodes}
samples with zero missing IDs, extras, duplicates, metadata mismatches, or
numeric trigger-code count mismatches. \\
\hline
\texttt{artifacts/validation\_runs/one\_block\_actual\_condition\_current}
& Accepted actual-condition one-block runner session and matching direct
loopback comparison. The current accepted session is
\texttt{P001\_20260612\_182812}. \\
\hline
\texttt{artifacts/validation\_runs/loopback\_calibration\_*}
& Internal package calibration artifacts used to verify the safe playback
workflow and report format. The accepted electrical skew estimate in this
report comes from the channel-scaled three-channel dummy direct-loopback run. \\
\hline
\texttt{artifacts/validation\_runs/dummy\_pulse\_*}
& Internal coded three-channel dummy WAV captures used to verify the analyzer
and acquisition workflow. The accepted functional run is
\texttt{final\_functional\_dummy\_pulse\_20260612\_channel\_scaled\_lslprobe}. \\
\hline
\texttt{artifacts/validation\_runs/audio\_route\_stress\_20260612\_131802} &
Silent callback stress showing Komplete ASIO device 31 supports stable
3-channel output while stereo Windows endpoints do not. \\
\hline
\texttt{artifacts/validation\_runs/mouse\_response\_timing\_lsl\_monotonic\_20260612\_134927}
with paired rich/numeric LSL probe folders & Internal mouse/response timing
stress after monotonic QC fix. Local logs contained 50 mouse clicks and 50
linked response markers; rich and numeric external LSL probes captured 104/104
samples each. \\
\hline
\texttt{artifacts/validation\_runs/mouse\_response\_timing\_lsl\_monotonic\_20260612\_134927/lsl\_reconciliation}
& Reconciliation report comparing local \texttt{events.csv},
\texttt{lsl\_markers.csv}, rich \texttt{PPSMarkersV2} probe rows, and numeric
\texttt{PPSTriggerCodes} probe rows. It passed with 104 compared events, no
missing/extra/duplicate event IDs, no metadata field mismatches, and matching
numeric trigger-code counts. \\
\hline
\texttt{artifacts/validation\_runs/mouse\_response\_timing\_lsl\_monotonic\_20260612\_134927/response\_strategy\_comparison}
& Strategy comparison report quantifying local mouse timing, rich LSL sample
timestamps, external LSL arrival timing, and callback-derived response-marker
timestamps for 50 paired click/marker events. \\
\hline
\texttt{artifacts/validation\_runs/session\_runner\_click\_path\_stress\_20260612\_deferred\_lsl}
& Session-runner click-path stress after deferring mouse-click LSL publication
until after the tactile response marker trigger. It produced 25/25 linked
click/marker pairs with \texttt{dac\_time\_sample\_exact} marker timestamps. \\
\hline
\texttt{validation\_protocols/scripts/compare\_response\_marker\_loopback.py}
& Internal analyzer for response-marker recovery from tactile-channel loopback
recordings. It is ready for physical session recordings and covered by a
synthetic recovery test. \\
\hline
\texttt{validation\_protocols/reports/evidence\_audit/}
& Internal requirement-level completeness audit. It is a checklist for what has
been proven in the current electrical setup and what has been intentionally
deferred; it records the accepted functional-route electrical baseline and
response-marker physical recovery. \\
\hline
\texttt{docs/WINDOWS\_PC\_SOFTWARE\_REQUIREMENTS.md} and
\texttt{artifacts/validation\_runs/pc\_software\_requirements\_*}
& Tracked install checklist plus ignored machine-readable audits of Python
packages, external tools, LabRecorder availability, and Komplete ASIO driver
visibility. \\
\hline
\end{longtable}

\section{Measured Results}

\subsection{Functional Three-Channel Direct Loopback Status}

\begin{longtable}{|p{0.30\textwidth}|p{0.24\textwidth}|p{0.32\textwidth}|}
\hline
\textbf{Claim} & \textbf{Current status} & \textbf{Publication rule} \\
\hline
Komplete ASIO three-channel playback & Accepted from callback stress testing:
the device opens as one 3-channel ASIO stream at 44100 Hz, blocksize 256, and
requested latency 0.010 s. & This supports using one multichannel WAV as the
stimulus-administration source. \\
\hline
Three-channel electrical identity & Passed. The channel-scaled safe route sweep
accepted output 1/input 1, output 2/input 2, and output 3/input 3 with all
channels visible, accepted for baseline detection, and unclipped. & Report
channel identity from this accepted direct-loopback route sweep. \\
\hline
Left/right electrical skew & Passed. Simultaneous dummy direct loopback:
0.023 ms mean absolute skew, SD 0.000 ms, paired count 5. & Criterion is at
most 1 ms. \\
\hline
Audio/tactile electrical skew & Passed. Tactile-vs-audio electrical skew:
0.011 ms mean absolute skew, SD 0.000 ms, paired count 5. & Criterion is at
most 2 ms. \\
\hline
Residual jitter and clipping/signal QC & Passed. Detection was 5/5 per channel,
all channel detection rates were 1.000, all low-signal flags were false, and
all clipping flags were false. Channel mean latencies were 33.605, 33.583, and
33.605 ms. & Accepted baseline shows at least 95\% pulse detection per channel,
p95 residual jitter at most 2 ms, maximum residual jitter at most 5 ms, no
clipping, and no low-signal flags. \\
\hline
\end{longtable}

Earlier route and level-adjustment captures are retained locally as setup
records. The defensible electrical latency result is the channel-scaled
functional route state recorded under
\texttt{artifacts/validation\_runs/final\_functional\_dummy\_pulse\_20260612\_channel\_scaled\_lslprobe/}.

\subsection{One-Block Capture Dry Run}

After the baseline was accepted, a single dummy block was rerun through the
same Komplete ASIO three-channel route and captured through direct electrical
loopback. The purpose was practical: verify that a whole block acquisition
lands immediately in analysis-ready files rather than only in a raw recording.
The run is stored under
\texttt{artifacts/validation\_runs/one\_block\_trial\_capture\_current/}.

\begin{longtable}{|p{0.32\textwidth}|p{0.26\textwidth}|p{0.30\textwidth}|}
\hline
\textbf{Output} & \textbf{Artifact} & \textbf{Analysis use} \\
\hline
Planned block timing & \texttt{planned\_pulses.csv/json} & Expected pulse
sample indices by event and channel. \\
\hline
Administered stimulus & \texttt{dummy\_3ch\_pulses.wav} & The one-block
three-channel WAV sent to the Komplete ASIO output. \\
\hline
Physical capture & \texttt{direct\_loopback\_capture.wav} & Raw electrical
record of outputs 1/2/3 returning to inputs 1/2/3. \\
\hline
Detected events & \texttt{dummy\_pulse\_detections.csv} & Row-level detected
sample index, latency, residual, source-channel identity, and shape-correlation
values. \\
\hline
Summary reports & \texttt{dummy\_pulse\_comparison\_report.json/md} & Immediate
pass/fail, per-channel latency distributions, inter-channel skew, clipping, and
low-signal checks. \\
\hline
Marker mirror & \texttt{lsl\_emitted\_markers.csv} & Software marker rows for
the five planned dummy pulse events. \\
\hline
\end{longtable}

The one-block dry run passed. It recovered 5/5 pulses on all three channels
with no clipping and no low-signal flags. Channel latency means were 33.465 ms
on channel 1, 33.447 ms on channel 2, and 33.469 ms on channel 3. The
left/right paired skew was 0.018 $\pm$ 0.010 ms, and the tactile-vs-audio
paired skew was 0.014 $\pm$ 0.005 ms. The generated CSV and JSON reports are
therefore already in the format needed for immediate analysis of channel
identity, latency, jitter, and skew.

\subsection{Actual-Condition One-Block Recording-Layer Run}

The accepted actual-condition recording-layer run is stored under
\texttt{artifacts/validation\_runs/recording\_layer\_actual\_condition\_20260613\_anchor/};
the accepted session folder is
\texttt{sessions/P001\_20260613\_174019/}. This run used the real prepared
Pfeiffer-style Segment 6 manifest, materialized one 88.0 s participant block,
played it through the Komplete ASIO endpoint, injected deterministic simulated
clicks after tactile onsets, wrote the optional local audio evidence WAV, wrote
the callback-derived local event/LSL marker mirrors, and recorded a
validation-only full-duplex direct electrical loopback trace. Conservative
validation gains were used (\texttt{audio\_gain = 0.0005},
\texttt{tactile\_gain = 0.02}) to keep the Komplete Audio 6 route out of
clipping while preserving timing evidence.

\begin{longtable}{|p{0.34\textwidth}|p{0.22\textwidth}|p{0.30\textwidth}|}
\hline
\textbf{Measure} & \textbf{Observed} & \textbf{Decision} \\
\hline
Actual-condition audit & passed & exactly one actual
\texttt{participant\_block\_wavs} block, no dummy/fake source flags \\
\hline
Scheduled callback-derived markers & 22 trial starts, 12 looming onsets, 20
tactile onsets, 22 response-window onsets, 22 trial ends & counts match the
prepared block CSV; all scheduled events were sample-exact \\
\hline
Immediate analysis output & 20 rows in
\texttt{analysis\_ready\_trials.csv} & one response row per tactile-onset
trial; catch trials remain catch trials \\
\hline
Internal XDF and LSL mirror & 146 loadable XDF samples; 146
\texttt{lsl\_markers.csv} rows & local/XDF/LSL mirror outputs are complete for
the one-block run; marker rows are emitted from runtime callback events, not
reconstructed from the planned CSV after playback \\
\hline
Local audio evidence WAV & 88.038 s, 3 channels, 3,882,496 frames, zero dropped
buffers, no clipped channels & data-heavy local software copy of the exact
output buffers generated by the runner \\
\hline
Direct electrical capture & 3 channels, no clipping & validation-only physical
reference for output arrival and channel synchronization \\
\hline
Response timing & RT mean 241.623 ms, SD 28.244 ms; mouse-to-marker delay mean
21.368 ms, SD 4.277 ms & simulated-click output is directly analyzable; marker
delay reflects the persistent full-duplex callback path in this validation run
\\
\hline
\end{longtable}

The recording-layer comparison is stored under
\texttt{artifacts/validation\_runs/recording\_layer\_alignment\_20260613\_anchor/}.
It compares the physical loopback WAV, local digital audio evidence WAV,
\texttt{events.csv}, and \texttt{lsl\_markers.csv} against the same sample
timeline. Physical-minus-digital latency was 33.462 $\pm$ 0.013 ms
(median 33.469 ms, p95 33.469 ms, min 33.447 ms, max 33.469 ms; channel count
3). Inter-channel synchronization remained well inside the acceptance
thresholds: right-minus-left was -0.023 ms, and tactile-minus-audio-mean was
0.011 ms. The internal marker mirror contained 146/146 event IDs with no
missing, extra, or duplicate IDs. Across 118 sample-indexed marker rows, the
callback-derived LSL timestamp error relative to the anchored sample timeline
was 0.00000019 $\pm$ 0.00000035 ms, median 0.000000 ms, p95 0.00000099 ms.
The physical tactile-channel response-marker trace recovered 20/20 expected
markers; after fitting the common recording offset, residual jitter was
0.000 ms mean, SD 0.000 ms, p95 0.000 ms.

\subsection{Audio-Device Stress}

\begin{longtable}{|p{0.34\textwidth}|p{0.22\textwidth}|p{0.30\textwidth}|}
\hline
\textbf{Measure} & \textbf{Observed} & \textbf{Decision} \\
\hline
Komplete ASIO 3-channel callback output & 3/3 iterations passed for 10 s each,
44100 Hz, blocksize 256, requested latency 0.010 s & pass \\
\hline
Actual reported output latency & 0.016508 s for the passing ASIO callback
stress rows & usable current default \\
\hline
Callback status flags & 0 status flags in all three passing ASIO rows & pass \\
\hline
Windows MME/DirectSound/WASAPI Komplete endpoints & all stereo endpoints failed
3-channel open checks because they expose only 2 output channels & expected;
do not use for 3-channel stimulus playback \\
\hline
\end{longtable}

\subsection{LSL and Local Log Reliability}

\begin{longtable}{|p{0.34\textwidth}|p{0.18\textwidth}|p{0.34\textwidth}|}
\hline
\textbf{Measure} & \textbf{Observed} & \textbf{Decision rule} \\
\hline
\texttt{PPSMarkersV2} resolved & yes & Exactly one intended rich stream
resolves for the run. \\
\hline
Probe event count & 12/12 for rich v2 burst; 5/5 for dummy pulse marker stream
& Matches expected actual marker count or documented dummy pulse count. \\
\hline
Dual-stream probe count & paced simultaneous test received 8/8 rich
\texttt{PPSMarkersV2} samples and 8/8 numeric \texttt{PPSTriggerCodes} samples
& Both streams should be available to external recorders during the same run. \\
\hline
LabRecorder XDF count & 16/16 rich \texttt{PPSMarkersV2} samples and 16/16
numeric \texttt{PPSTriggerCodes} samples; zero missing, extra, duplicate, or
field-mismatched rich event IDs; numeric code counts matched & External XDF
files should preserve the same marker set as the local LSL mirror. \\
\hline
Missing LSL event IDs & zero in tested bursts & Must be zero. \\
\hline
Duplicate LSL event IDs & zero in tested bursts & Must be zero. \\
\hline
Trigger-code consistency & v2 probe preserved event codes and trigger keys;
local \texttt{trigger\_dictionary.json} was written for the burst & Codes match
\texttt{trigger\_dictionary.json}; full participant runner validation still
needed. \\
\hline
Mouse/response reconciliation & 104/104 local actual events matched 104/104
rich LSL samples; 104/104 numeric trigger-code samples matched expected code
counts; zero field mismatches & Rich LSL plus local logs should reconstruct the
same actual event set. Numeric LSL should match trigger-code counts, but cannot
recover event identity alone. \\
\hline
Response strategy timing & Local marker-minus-mouse mean 8.000 ms, SD 0.000
ms, median 8.000 ms; rich LSL mouse-sample minus local mouse mean 0.081 ms, SD
0.056 ms, median 0.067 ms; rich LSL mouse-arrival minus local mouse mean 0.421
ms, SD 0.223 ms, median 0.339 ms; response-marker LSL arrival minus marker
timestamp median -5.105 ms & Local mouse logs are the primary response-time
source; rich LSL sample timestamps are a close external record; LSL arrival
time is a monitoring metric, not the event timestamp. \\
\hline
Session runner click path & 25/25 controller-level clicks generated 25/25
linked response markers during active playback; marker-minus-mouse mean 8.150
ms, SD 0.081 ms, median 8.105 ms, max 8.317 ms, all marker timestamps
\texttt{dac\_time\_sample\_exact} & The session runner can link clicks to
callback-derived response markers without putting LSL publishing before the
physical marker trigger. \\
\hline
Visible runner OS click path & 10/10 armed OS clicks into the retired Tk runner
target were logged as in-target clicks during active fake-audio playback; 10/10
linked response markers were generated; marker-minus-mouse mean 8.553 ms, SD
0.125 ms, median 8.541 ms, p95 8.762 ms & Preserves historical visible-runner
click delivery evidence without touching the audio interface. Focus Mode is the
current operational path. \\
\hline
One-block realtime runner output path & 5/5 tactile onsets, 5/5 mouse clicks,
5/5 response markers, 5/5 trial-end markers, no duplicate event IDs; pyxdf
loaded 43 internal-XDF samples; \texttt{analysis\_ready\_trials.csv} contained
5/5 hits; simulated RT mean 223.215 ms, SD 30.936 ms; mouse-to-marker mean
8.081 ms, SD 0.014 ms & Demonstrates that a single prepared block can produce
immediately analyzable local CSV/XDF/marker artifacts under realtime software
conditions. This is not a physical loopback or hardware-latency estimate. \\
\hline
Timestamp-quality distribution & 7 \texttt{software\_log}; 5
\texttt{dac\_time\_sample\_exact} in the v2 burst & Expected callback-derived
qualities dominate for scheduled audio events; fallbacks require review. \\
\hline
Arrival-minus-marker behavior & software-log v2 markers arrived about
0.15--0.45 ms after their marker timestamps; dummy pulse markers arrived about
0.20--0.35 ms after their marker timestamps. Synthetic future callback markers
showed negative arrival-minus-marker values because the burst emitted future
timestamps immediately. & Report median, p95, and max; interpret as recording
behavior only. \\
\hline
Outlet lifetime behavior & an unpaced dual-stream burst resolved both streams
but captured 0/8 samples on both probes; a paced version with outlet warmup
captured 8/8 on both streams & Create LSL outlets before the first important
marker, keep them alive for the run, and avoid one-shot marker streams. \\
\hline
\end{longtable}

\subsection{Recording-Layer Reliability Comparison}

\begin{longtable}{|p{0.22\textwidth}|p{0.25\textwidth}|p{0.29\textwidth}|p{0.14\textwidth}|}
\hline
\textbf{Recording layer} & \textbf{Reliability evidence} & \textbf{Timing interpretation} & \textbf{Current role} \\
\hline
Local event log and internal marker mirror & The actual-condition block wrote
146 runtime event IDs to \texttt{events.csv} and 146 matching rows to
\texttt{lsl\_markers.csv}, with no missing, extra, or duplicate marker IDs.
Earlier stress runs also reconciled 104/104 mouse/response events. & Primary
software record. Mouse-click time is taken from the local process clock;
scheduled stimulus markers are emitted when their sample positions cross the
audio callback timeline. & Primary analysis and reconstruction source. \\
\hline
LSL rich/numeric streams and LabRecorder/XDF & External probes recovered 5/5
dummy pulse markers, 12/12 rich v2 burst markers, 8/8 paced rich/numeric
dual-stream markers, 104/104 mouse/response actual events, and LabRecorder/XDF
captured 16/16 rich plus 16/16 numeric samples with zero field mismatches. Rich
LSL mouse sample lag was 0.081 ms mean, SD 0.056 ms; probe arrival lag was
0.421 ms mean, SD 0.223 ms. In the actual-condition internal marker mirror,
sample-anchored LSL timestamp error had p95 0.00000099 ms. & LSL sample
timestamps are suitable for external software reconstruction. Arrival time
measures recorder/client behavior and is not a physical onset timestamp. &
External reconstruction and synchronization with other acquisition systems. \\
\hline
Local audio evidence WAV & The actual-condition block wrote a 3-channel,
88.038 s digital output evidence WAV with zero dropped buffers and no clipped
channels. Physical-minus-digital offset was 33.462 $\pm$ 0.013 ms when compared
with direct loopback. & This is a local software copy of the generated output
buffers after routing, gain, clipping limits, and response-marker injection; it
is not a physical arrival measurement. & Optional data-heavy safety copy during
participant runs. \\
\hline
WASAPI loopback & WASAPI was requested during ASIO multichannel dummy playback
but produced no usable capture in the tested configuration. & It records the
Windows endpoint, not necessarily the Native Instruments ASIO stream, so it
does not observe the physical three-channel route under test. & Excluded from
the core validation strategy. \\
\hline
Direct electrical physical loopback & Channel-scaled route sweep passed
identity and signal QC. Simultaneous dummy direct loopback recovered 5/5 pulses
on all three channels with no clipping or low-signal flags. Per-channel latency
means were 33.605, 33.583, and 33.605 ms; left/right skew was
0.023 $\pm$ 0.000 ms; tactile/audio skew was 0.011 $\pm$ 0.000 ms. & This is
the ground-truth layer for electrical inter-channel latency and skew in the
current setup. Woojer mechanical vibration onset is a later hardware phase. &
Accepted electrical timing baseline. \\
\hline
\end{longtable}

\section{Belts-and-Suspenders Data Strategy}

The validation supports a layered recording strategy rather than relying on a
single timing record. The recommended strategy is:

\begin{enumerate}
  \item Use the callback-derived local event log as the primary software record.
  Scheduled stimulus events should be timestamped when their sample index enters
  the active output callback buffer, preferably from the output DAC time when
  PortAudio exposes it.
  \item Publish the same actual events to \texttt{PPSMarkersV2}. This rich LSL
  stream is the external, multimodal synchronization record and should be
  sufficient to reconstruct the experiment if captured by LabRecorder or another
  LSL client.
  \item Publish the numeric \texttt{PPSTriggerCodes} stream at the same time for
  systems that can only record integer triggers. Reconstruction from this stream
  depends on the session's \texttt{trigger\_dictionary.json}.
  \item Keep local \texttt{events.csv}, \texttt{lsl\_markers.csv}, and
  \texttt{trigger\_dictionary.json} as the non-network backup record.
  \item Enable the local audio evidence WAV when the operator wants a data-heavy
  safety copy of the generated output buffers. This WAV is useful for recovering
  the exact generated multichannel stimulus stream and tactile-channel
  response-marker clicks if a marker recorder fails.
  \item Keep physical loopback as a validation and publication-calibration
  layer, not as normal participant-run equipment. Direct electrical loopback is
  the correct way to validate channel timing and to compare the event/LSL and
  local audio evidence layers against physical arrival. WASAPI loopback was
  evaluated and removed from the core strategy for this ASIO route because it
  did not capture the multichannel playback stream.
  \item Feed mouse-click response markers into the tactile/audio recording path
  only as a backup/reconstruction trace. The immediate software mouse-click log
  remains the primary response event unless a future button box or hardware
  trigger device is added.
  \item For mouse responses, trigger the tactile response marker before pushing
  the mouse event to LSL. The mouse event is still logged locally first so the
  response marker can carry the mouse event ID, but LSL/network work should not
  sit in front of the physical backup marker.
\end{enumerate}

In plain language, this means the experiment should write down what it actually
did, broadcast those same runtime events through LSL, and optionally save the
generated output buffers locally when a second software safety copy is desired.
Physical loopback is the independent validation witness used to prove those
software layers are aligned; it is not required for ordinary participant runs.

\subsection{Mouse Click and Response Marker}

\begin{longtable}{|p{0.34\textwidth}|p{0.18\textwidth}|p{0.34\textwidth}|}
\hline
\textbf{Measure} & \textbf{Observed} & \textbf{Decision rule} \\
\hline
Emulated OS click count & 10/10 armed OS clicks were delivered to the visible
legacy runner target during active fake-audio playback and logged in target. &
Every armed click during playback is logged. \\
\hline
Linked response-marker count & Synthetic timing stress linked 50/50 clicks;
session-controller stress linked 25/25 clicks; visible OS-click stress linked
10/10 clicks to \texttt{response\_marker\_start}. & Every in-playback click has
a \texttt{response\_marker\_start}. \\
\hline
Session-runner click-path linkage & 25/25 runner-level clicks linked to 25/25
response markers with \texttt{dac\_time\_sample\_exact} timestamps & Validates
the participant session controller path separately from OS desktop-click
delivery and physical loopback. \\
\hline
Response-marker loopback analyzer & \implemented{} script with synthetic
validation & Compares logged \texttt{response\_marker\_start} sample indices
against tactile-channel loopback pulses, fits the hardware/recording offset,
and reports detection rate plus residual jitter. \\
\hline
Physical response-marker recovery & 5/5 response-marker-style channel-3 pulses
were recovered from the accepted physical direct-loopback recording. Fitted
offset was 33.605 ms mean, SD 0.000 ms; absolute residual jitter was
0.000 ms mean, SD 0.000 ms. & Report detection rate, fitted offset, p95
residual, and maximum residual. \\
\hline
Local-vs-LSL mouse timing & local click-to-marker delay is exact in the
software harness; rich LSL mouse sample timestamps lag local click times by
about 0.067 ms median, and probe arrivals by about 0.339 ms median & Use local
mouse log as the primary response timestamp; use LSL sample timestamps for
external reconstruction and arrival time for health monitoring. \\
\hline
Physical tactile-channel marker recovery & \proven{} electrically for the
channel-3 loopback path; Woojer mechanical onset is a later hardware phase. &
Loopback recording contains expected marker pulses. \\
\hline
\end{longtable}

\section{Repeatable Electrical Validation Run}

The electrical three-channel baseline is now accepted for the tested functional
route. For reproducibility, the accepted electrical route can be rerun with the same
channel-scaled validation amplitudes:

\begin{verbatim}
python .\For-AI/engineering/validation\scripts\run_dummy_output_route_sweep.py ^
  --device-query Komplete --amplitude 0.02 ^
  --channel-amplitudes 1:0.0005,2:0.02,3:0.02 ^
  --input-channels 6 --output-channels 3 --sweep-output-count 3

python .\For-AI/engineering/validation\scripts\run_dummy_pulse_latency.py ^
  --device-query Komplete --amplitude 0.02 ^
  --channel-amplitudes 1:0.0005,2:0.02,3:0.02 ^
  --record-asio-loopback --emit-lsl ^
  --input-channels 3 --output-channels 3

python .\For-AI/engineering/validation\scripts\compare_dummy_pulse_recordings.py ^
  --run-dir artifacts\validation_runs\dummy_pulse_YYYYMMDD_HHMMSS
\end{verbatim}

The validation scripts refuse direct-loopback hardware playback above amplitude
0.10. If future signal levels are unacceptable, adjust physical gain, input
mode, or routing rather than raising the digital playback level beyond that
safety cap.

Woojer mechanical vibration onset is intentionally outside this current run
because the Woojer is not physically in the validation loop. When that hardware
is added, the added measurement should record an external vibration/contact
sensor alongside the channel-3 electrical drive.

\section{Progress Log}

\begin{longtable}{|p{0.18\textwidth}|p{0.72\textwidth}|}
\hline
\textbf{Date} & \textbf{Progress} \\
\hline
2026-06-12 & Created internal validation protocol folder outside the runtime
package. \\
\hline
2026-06-12 & Added dummy three-channel pulse stimulus generator, direct
loopback runner, comparison script, acceptance criteria, and synthetic tests. \\
\hline
2026-06-12 & Updated validation protocol documentation to make single-file
channel routing, audio/tactile electrical skew, and measurement-pipeline
latency explicit. \\
\hline
2026-06-12 & Updated LSL probe/protocol to target \texttt{PPSMarkersV2} and
normalize rich marker rows for comparison against local event logs. \\
\hline
2026-06-12 & Added this living LaTeX report so hardware results can be entered
without erasing unmeasured limitations. \\
\hline
2026-06-12 & Installed the optional LSL dependency and validated both the dummy
pulse marker stream and the rich \texttt{PPSMarkersV2} stream with external
probes. Dummy markers were received 5/5; the rich v2 burst was received 12/12
with no duplicate event IDs. \\
\hline
2026-06-12 & Ran preliminary package loopback calibration and audio-device
stress. The Komplete ASIO endpoint passed 3-channel callback stress. Later
channel-scaled dummy direct-loopback testing supplied the accepted all-channel
electrical skew baseline. \\
\hline
2026-06-12 & Ran simultaneous rich/numeric LSL probing. The first unpaced
one-shot burst resolved both streams but captured no samples, showing that very
short LSL outlet lifetimes are fragile. A paced run with outlet warmup captured
8/8 \texttt{PPSMarkersV2} samples and 8/8 \texttt{PPSTriggerCodes} samples. \\
\hline
2026-06-12 & Added and ran a single-output route identity sweep. Route identity
and level checks are retained as internal setup records. The later
channel-scaled functional sweep accepted all three physical routes without
clipping or low-signal flags. \\
\hline
2026-06-12 & Added hardware-playback safety guards. Validation scripts now
default to amplitude 0.05 and refuse amplitudes above 0.10 so weak channels are
handled by gain/routing adjustment rather than unsafe playback levels. \\
\hline
2026-06-12 & Added an ASIO output-selector sweep option and clarified the gain
criterion. Equal physical knob positions are not sufficient evidence; route
acceptance depends on recorded peak, SNR, shape correlation, and no clipping.
The channel-scaled functional sweep now passes that shared signal-quality gate. \\
\hline
2026-06-12 & Added per-channel dummy validation amplitudes so the internal
validation WAV can avoid clipping on high-gain capture paths while keeping
lower-gain paths measurable. The accepted settings were channel 1 = 0.0005,
channel 2 = 0.02, and channel 3 = 0.02. The accepted simultaneous direct
loopback recovered 5/5 pulses on all channels with no clipping, no low-signal
flags, left/right skew 0.023 $\pm$ 0.000 ms, and tactile/audio skew
0.011 $\pm$ 0.000 ms. An external LSL probe recorded 5/5 dummy markers; WASAPI
initialized but recorded no data and was removed from the core ASIO validation
strategy. \\
\hline
2026-06-12 & Added and ran the internal mouse/response timing stress harness.
The monotonic timing run logged 50/50 clicks and 50/50 linked
response markers, while external rich and numeric LSL probes captured 104/104
samples on each stream. \\
\hline
2026-06-12 & Added and ran LSL/local-log reconciliation for the mouse/response
stress run. The rich stream matched all 104 local actual event IDs and marker
metadata with zero mismatches; the numeric stream matched all expected
trigger-code counts. \\
\hline
2026-06-12 & Added and ran response timing strategy comparison. In the
50-click stress run, rich LSL mouse-click sample timestamps followed local
mouse logs by about 0.067 ms median, while external LSL probe arrival followed
by about 0.339 ms median. Response-marker sample timestamps remained aligned to
the planned 8 ms callback-derived marker delay; LSL arrival time is documented
as a monitoring metric, not the event timestamp. \\
\hline
2026-06-12 & Ran safe-amplitude direct-loopback setup checks and kept them as
internal acquisition records. Later channel-scaled dummy direct-loopback testing
provided the accepted electrical inter-channel skew values after passing the
shared detection, jitter, clipping, and signal-quality criteria.
\\
\hline
2026-06-12 & Hardened \texttt{compare\_dummy\_pulse\_recordings.py} so pulse
detection is tied to local windows around planned sample positions. This
prevents startup artifacts or missing channels from producing misleading skew
values. \\
\hline
2026-06-12 & Added a plain-language report section explaining why the dummy
three-channel pulse stimulus was used, what has been accepted, and what the
measurement boundaries are. \\
\hline
2026-06-12 & Added a requirement-level evidence audit under
\texttt{validation\_protocols/reports/evidence\_audit/}. The audit is internal
and conservative; the LaTeX report remains the publication-facing validation
narrative. \\
\hline
2026-06-12 & Added and ran
\texttt{run\_session\_runner\_click\_path\_stress.py}. The current runner
logs mouse clicks locally, triggers the tactile response marker, then pushes
the mouse marker to LSL. The stress logged 25/25 linked click/marker pairs with
marker-minus-mouse delay mean 8.150 ms, SD 0.081 ms, median 8.105 ms, and
maximum 8.317 ms for an intended 8 ms marker. \\
\hline
2026-06-12 & Added
\texttt{compare\_response\_marker\_loopback.py}, an internal checker that
compares logged \texttt{response\_marker\_start} sample indices with
tactile-channel loopback pulses. Synthetic validation verifies that it recovers
known marker pulses, fitted offset, and residual jitter; the accepted
channel-3 dummy loopback was also recovered as a response-marker-style physical
trace. \\
\hline
2026-06-12 & Added the Windows PC software requirements checklist and a
machine-readable dependency audit. Installed \texttt{pyxdf} for offline XDF
inspection and downloaded/extracted the official LabRecorder Windows release
under ignored local folders. The audit now reports no missing runtime,
validation, or external-tool dependencies on this PC. \\
\hline
2026-06-12 & Added and ran the LabRecorder/XDF stress harness. The current run
captured 16/16 rich \texttt{PPSMarkersV2} samples and 16/16 numeric
\texttt{PPSTriggerCodes} samples in XDF with no missing, extra, duplicate, or
field-mismatched rich events and no numeric trigger-code count mismatches. \\
\hline
2026-06-12 & Added and ran
\texttt{run\_visible\_runner\_os\_click\_stress.py}. The harness opened the
retired visible Tk runner with a deterministic fake audio engine and sent armed
OS clicks into the click target. The first run exposed a Windows
DPI-scaling coordinate mismatch in the click-target hit test; after correcting
the runner to use DPI-adjusted bounds, the stress recovered 10/10 in-target
clicks and 10/10 linked response markers during playback. \\
\hline
2026-06-12 & Updated validation analysis outputs and this report to require
scientific timing summaries as mean $\pm$ SD, with paired inter-channel skew
statistics for left/right and tactile-vs-audio synchronization. The final
report must compare local logs, LSL/XDF, and direct physical loopback
reliability as separate recording layers, while documenting that WASAPI was
evaluated and excluded for this ASIO route. \\
\hline
2026-06-12 & Ran a single dummy block through the Komplete ASIO three-channel
route and direct electrical loopback after the baseline decision. The run wrote
the source WAV, planned pulse CSV/JSON, raw loopback WAV, marker mirror,
row-level detection CSV, and JSON/Markdown analysis reports under
\texttt{artifacts/validation\_runs/one\_block\_trial\_capture\_current/}. It
passed with 5/5 pulses detected on all channels, left/right paired skew
0.018 $\pm$ 0.010 ms, and tactile-vs-audio paired skew
0.014 $\pm$ 0.005 ms. \\
\hline
2026-06-13 & Added the optional local audio evidence WAV path to the runner
capture stack and the recording-layer comparison protocol. The audio evidence
WAV is written from the mixed output callback and is compared against
validation-only physical loopback, while \texttt{events.csv},
\texttt{lsl\_markers.csv}, and local XDF files remain the lightweight marker
record. \\
\hline
2026-06-13 & Updated callback-derived scheduled marker timestamps to use the
\texttt{audio\_sample\_zero} sample-time anchor within each block. This removes
per-buffer PortAudio timing drift from later scheduled trial markers while
preserving DAC-time quality labels when the anchor is available. \\
\hline
2026-06-13 & Ran one actual prepared Segment 5/6 block with physical loopback,
local digital audio evidence WAV, and internal event/LSL marker mirror enabled.
The recording-layer alignment passed with physical-minus-digital latency
33.462 $\pm$ 0.013 ms, right-minus-left skew -0.023 ms, tactile-minus-audio
skew 0.011 ms, zero dropped digital buffers, no clipping, 146/146 internal
marker IDs, p95 sample-anchored LSL timestamp error below 0.000001 ms, and
20/20 physical response-marker pulses recovered with 0.000 ms residual jitter
after fitting the common recording offset. \\
\hline
2026-06-13 & Added the main conclusion section with plain-language and technical
interpretation of LSL/event records, local audio evidence WAVs, and physical
loopback. The conclusion now includes accepted latency/reliability trial tables
and a recording-method assessment for the final validation report. \\
\hline
\end{longtable}

\section{Main Conclusion}

\subsection{Plain-Language Conclusion}

The validation supports the current recording strategy. The experiment runner
can administer one synchronized three-channel audio-tactile block through the
Komplete Audio 6, record the runtime events needed for analysis, and optionally
write a local audio evidence WAV that matches the physical loopback after one
stable hardware delay is accounted for. The physical loopback is later than the
LSL/event record because it measures the electrical signal after the audio
interface output, cable path, and input capture stage. This delay is not an LSL
failure. It is the expected physical delay between a software-scheduled audio
sample and the same signal appearing in the validation loopback recording.

The key comparison is whether the timing information recovered from the WAV
recordings agrees with the timing information directly logged by the runner and
mirrored to LSL. For the events that are visible in the audio/tactile waveform,
the answer is yes. In the accepted actual-condition block, the physical loopback
was a stable delayed copy of the runner's local audio evidence WAV. The fitted
physical-minus-digital delay was 33.462 ms with an SD of 0.013 ms. The physical
tactile-channel response-marker pulses were recovered 20/20 times, and after
the common physical offset was fitted, the residual marker error was 0.000 ms.
Thus, the WAV-derived evidence and the LSL/event-derived evidence are
consistent where their information overlaps.

The recording layers serve different purposes. The local event log and internal
LSL marker mirror are the clean structured record: they contain the participant,
block, trial type, trigger code, sample index, tactile cue, mouse click, response
marker, and response outcome information needed for immediate analysis. The
local audio evidence WAV is a data-heavy backup of the generated output buffers;
it can show that the runner produced the expected multichannel waveform and
response-marker clicks. The physical loopback is the validation reference that
proves the generated signals actually arrived at the electrical outputs with
the measured delay and channel skew. In normal participant experiments, the
recommended configuration is therefore local event logs plus internal LSL marker
mirrors by default, with the optional local audio evidence WAV enabled when a
second software safety copy is desired. Physical loopback is reserved for
validation and publication calibration, not routine data collection.

The channel differences are also consistent across the waveform-based evidence.
The direct dummy baseline and the actual-condition recording-layer comparison
both show sub-millisecond inter-channel differences, on the order of one audio
sample. In the accepted actual-condition recording-layer comparison, the right
channel was 0.023 ms earlier than the left channel, and the tactile channel was
0.011 ms later than the mean of the two audio channels. These values are far
below the acceptance thresholds of 1 ms for left/right skew and 2 ms for
tactile-vs-audio skew. No evidence from the accepted runs suggests a
method-specific timing contradiction between the LSL/event record, the local
audio evidence WAV, and the physical loopback. LSL does not record the audio
waveform itself, so channel skew is not estimated from LSL alone; it is derived
by comparing LSL/event sample indices and WAV-detected signal features on the
same sample timeline.

\subsection{Technical Conclusion}

The accepted timing model is:

\begin{verbatim}
runtime sample event timestamp
  + stable physical output/input/capture delay
  = physical loopback detection time
\end{verbatim}

In the accepted actual-condition recording-layer run, the fitted physical
delay between the runner's digital audio evidence WAV and the physical
electrical loopback WAV was 33.462 $\pm$ 0.013 ms across the three channels
(median 33.469 ms, p95 33.469 ms, min 33.447 ms, max 33.469 ms). Internal
event/LSL marker timestamps were emitted from runtime callback events and then
anchored to the block's \texttt{audio\_sample\_zero} sample-time reference. The
sample-anchored LSL timestamp error relative to the runner's audio-sample
timeline had mean 0.00000019 ms, SD 0.00000035 ms, median 0.000000 ms, and p95
0.00000099 ms across 118 sample-indexed markers. This validates the internal
marker mirror as a sample-timeline record of runner timing. External
LabRecorder/XDF preservation was validated separately: 16/16 rich
\texttt{PPSMarkersV2} samples and 16/16 numeric \texttt{PPSTriggerCodes}
samples were preserved with zero missing, extra, duplicate, or field-mismatched
rich events.

The practical implication is that the final analysis should use the
callback-derived event/LSL layer for trial reconstruction and behavioral
outcomes, while using the local audio evidence WAV as an optional signal-level
backup. Physical loopback remains the reference method for validating absolute
electrical latency and inter-channel skew. WASAPI loopback is not part of the
accepted timing method for this setup because the Native Instruments ASIO route
can bypass the Windows endpoint that WASAPI records.

\subsection{Accepted Latency and Reliability Trials}

\begin{longtable}{|p{0.20\textwidth}|p{0.19\textwidth}|p{0.29\textwidth}|p{0.22\textwidth}|}
\hline
\textbf{Validation trial} & \textbf{Evidence count} & \textbf{Timing result} & \textbf{Conclusion} \\
\hline
Komplete ASIO route stress & 3/3 ten-second 3-channel callback iterations &
Reported output latency was 16.508 ms for the stable ASIO callback-stress
rows; callback status flags were 0. & The Komplete ASIO endpoint can sustain
one synchronized 3-channel output stream. \\
\hline
Channel-scaled dummy direct loopback baseline & 5/5 pulses detected on each of
3 channels & Channel means were 33.605, 33.583, and 33.605 ms; all channel SDs
were 0.000 ms. Left/right skew was 0.023 $\pm$ 0.000 ms; tactile/audio skew was
0.011 $\pm$ 0.000 ms. & This is the accepted electrical baseline for
single-file three-channel routing and inter-channel synchronization. \\
\hline
Dummy one-block physical capture & 5/5 pulses detected on all channels &
Channel means were 33.465, 33.447, and 33.469 ms. Left/right skew was
0.018 $\pm$ 0.010 ms; tactile/audio skew was 0.014 $\pm$ 0.005 ms. & The
analysis pipeline can turn a block-style loopback WAV directly into latency and
skew tables. \\
\hline
Actual-condition one-block runner audit & 22/22 trial starts, 12/12 looming
onsets, 20/20 tactile onsets, 22/22 response-window onsets, 22/22 trial ends &
The runner wrote 146 loadable XDF samples, 146 LSL marker rows, and 20
analysis-ready response rows. & One prepared Segment 5/6 block produces the
expected analysis-ready event and marker outputs. \\
\hline
Actual-condition recording-layer alignment & Physical loopback, digital audio
evidence WAV, \texttt{events.csv}, and \texttt{lsl\_markers.csv} compared on
one sample timeline & Physical-minus-digital delay was 33.462 $\pm$ 0.013 ms
(median 33.469 ms, p95 33.469 ms, range 33.447--33.469 ms). Right-minus-left
skew was -0.023 ms; tactile/audio skew was 0.011 ms. & The generated audio WAV
and physical loopback agree after one stable physical offset. \\
\hline
Internal LSL/event marker mirror in actual-condition block & 146/146 event IDs
mirrored; 118 sample-indexed marker rows used for timestamp comparison & No
missing, extra, or duplicate marker IDs. Sample-anchored LSL timestamp p95 error
was 0.00000099 ms. & The internal LSL-style record is an accurate
sample-timeline representation of runtime events. \\
\hline
Physical response-marker recovery in actual-condition block & 20/20
tactile-channel response-marker pulses recovered & After fitting the common
recording offset, residual jitter was 0.000 ms mean, SD 0.000 ms, p95
0.000 ms. & WAV-derived response-marker timing matches the event/LSL record
where response markers are physically encoded. \\
\hline
External LSL probe and LabRecorder/XDF preservation & 5/5 dummy markers, 12/12
rich burst markers, 8/8 paced rich/numeric markers, 104/104 mouse-response
events, and 16/16 LabRecorder rich plus 16/16 LabRecorder numeric markers &
Rich LSL mouse sample lag was 0.081 $\pm$ 0.056 ms relative to local mouse
logging; probe arrival lag was 0.421 $\pm$ 0.223 ms. LabRecorder XDF timestamp
delta relative to local marker mirror was about -0.005 ms. & External LSL/XDF
preserves the structured marker record. Arrival time is monitoring metadata, not
physical signal arrival. \\
\hline
Mouse-click and response-marker software path & 50/50 synthetic clicks, 25/25
session-controller clicks, and 10/10 visible-runner OS clicks linked to response
markers & Session-controller marker-minus-mouse delay was 8.150 $\pm$ 0.081 ms,
median 8.105 ms, maximum 8.317 ms. & Mouse clicks are logged locally first and
linked to tactile-channel response markers without placing LSL publishing before
the marker trigger. \\
\hline
WASAPI diagnostic attempt & WASAPI initialized during ASIO dummy playback but
recorded no usable multichannel data & No accepted latency estimate is derived
from WASAPI for this route. & WASAPI is excluded from the core validation
strategy for Native Instruments ASIO playback. \\
\hline
\end{longtable}

\subsection{Recording Method Assessment}

\begin{longtable}{|p{0.20\textwidth}|p{0.24\textwidth}|p{0.30\textwidth}|p{0.16\textwidth}|}
\hline
\textbf{Recording method} & \textbf{Reliability assessment} & \textbf{What it proves} & \textbf{Recommended use} \\
\hline
Local \texttt{events.csv} & Very high in accepted tests; all actual-condition
runtime events were present with no duplicate event IDs. & The runner's
structured trial and response history. & Always on; primary analysis record. \\
\hline
Internal \texttt{lsl\_markers.csv}/local XDF mirror & Very high in accepted
tests; 146/146 runtime event IDs mirrored, and sample-anchored timestamp error
was below 0.001 microseconds at p95. & The same runtime marker set in an
LSL-style structure, including trigger codes and trial metadata. & Always on
unless deliberately disabled for a special diagnostic. \\
\hline
External LSL/LabRecorder & High for marker preservation; separately validated
with complete rich and numeric marker capture. & Whether other acquisition
systems can receive and preserve the structured marker stream. & Use for EEG,
external acquisition, and lab synchronization. \\
\hline
Local audio evidence WAV & Very high for generated-output evidence; the accepted
run had zero dropped buffers and no clipped channels. & The exact multichannel
buffers the runner attempted to send after routing, gain, clipping limits, and
response-marker injection. & Optional data-heavy backup during participant
runs. \\
\hline
Physical electrical loopback & Highest for electrical timing validation; it is
the only layer that proves physical signal arrival and channel skew. & Absolute
electrical latency, inter-channel skew, crosstalk/channel-loss checks, and
physical response-marker recovery. & Validation and publication calibration
only. \\
\hline
WASAPI loopback & Not reliable for this ASIO multichannel route. & No accepted
physical timing claim in this setup. & Excluded except for optional Windows
endpoint diagnostics. \\
\hline
Woojer mechanical sensing & Not part of the current electrical loop. & Would
measure mechanical vibration onset, not just channel-3 electrical drive. &
Future sensor/contact-mic validation. \\
\hline
\end{longtable}

\section{Interpretation Rules}

The direct electrical loopback is the primary ground truth for whether one
three-channel file is split correctly by the interface and whether the
headphone and tactile electrical signals remain synchronized. LSL and local
logs are data-recording systems; they cannot prove physical signal arrival by
themselves. WASAPI loopback was evaluated and removed from the core strategy
because it does not reliably observe the Native Instruments ASIO multichannel
stream. Woojer mechanical vibration onset belongs to a later hardware phase
because the Woojer is not physically in the current validation loop.

\end{document}
